%%%%%%%%%%%%%%%%%%%%%%%%%%%%%%%%%%%%%%%%%%%%%%%%%%%%%%%%%%%%%%%%%%%%%%%%%%%%%%%%
% Tacet: Reliable Timing Side-Channel Detection via Bayesian Inference
% USENIX Security 2026 Submission
%
% Target: USENIX Security 2026, Cycle 2
% Submission Deadline: Thursday, February 5, 2026, 11:59 PM AoE
%%%%%%%%%%%%%%%%%%%%%%%%%%%%%%%%%%%%%%%%%%%%%%%%%%%%%%%%%%%%%%%%%%%%%%%%%%%%%%%%

\documentclass[letterpaper,twocolumn,10pt]{article}
\usepackage{usenix}

% Additional packages for figures and math
\usepackage{tikz}
\usepackage{amsmath}
\usepackage{amssymb}
\usepackage{graphicx}
\usepackage{booktabs}
\usepackage{multirow}
\usepackage{subcaption}
\usepackage{xspace}

% For code listings
\usepackage{listings}
\lstset{
  basicstyle=\ttfamily\small,
  breaklines=true,
  frame=single,
  columns=fullflexible,
}

% For URLs
\usepackage{hyperref}
\hypersetup{hidelinks}

% Custom commands
\newcommand{\tacet}{\textsc{Tacet}\xspace}
\newcommand{\dudect}{\textsc{dudect}\xspace}
\newcommand{\tlsfuzzer}{\textsc{tlsfuzzer}\xspace}
\newcommand{\rtlf}{\textsc{RTLF}\xspace}
\newcommand{\silent}{\textsc{SILENT}\xspace}
\newcommand{\ie}{i.e.,\xspace}
\newcommand{\eg}{e.g.,\xspace}
\newcommand{\etal}{et~al.\xspace}
\newcommand{\prob}[1]{P(#1)}
\newcommand{\given}{\,|\,}

%-------------------------------------------------------------------------------
\begin{document}
%-------------------------------------------------------------------------------

% No date on title page
\date{}

% Title (14pt bold per USENIX requirements)
\title{\Large \bf Tacet: Reliable Timing Side-Channel Detection\\via Bayesian Inference}

% Authors - ANONYMIZED for submission
% Replace with actual author info for camera-ready
\author{
{\rm Anonymous}\\
Anonymous Institution
% Add more authors as needed:
% \and
% {\rm Second Author}\\
% Second Institution
}

\maketitle

%-------------------------------------------------------------------------------
\begin{abstract}
%-------------------------------------------------------------------------------
% ~150-200 words summarizing the key contributions
% Hook: existing tools give false confidence
% Solution: calibrated Bayesian verdicts
% Results: highlight key evaluation findings

Timing side-channel attacks remain a persistent threat to cryptographic implementations.
Existing detection tools like \dudect and \tlsfuzzer provide binary pass/fail verdicts,
but offer no indication when measurement noise prevents reliable conclusions---leading
developers to place false confidence in potentially vulnerable code.
%
We present \tacet, a timing side-channel detection library that provides
\emph{calibrated verdicts}: when evidence is insufficient, \tacet reports
Inconclusive rather than guessing, indicating when evidence is insufficient.
%
\tacet achieves this through Bayesian inference over quantile differences,
outputting interpretable posterior probabilities $\prob{\text{leak} > \theta \given \text{data}}$
rather than p-values.
Key innovations include block bootstrap covariance estimation to handle
autocorrelated measurements, quality gates that detect when the model
cannot reliably distinguish signal from noise, and explicit measurement
floor tracking.
%
Our evaluation on synthetic benchmarks and real cryptographic libraries demonstrates
that \tacet maintains calibrated verdicts across varying noise levels and
autocorrelation conditions where existing tools produce unreliable results.
\tacet is implemented in Rust with C/C++, Go, and JavaScript bindings, designed for
integration into continuous integration pipelines.

\end{abstract}

%===============================================================================
% SECTION 1: INTRODUCTION (~1.25 pages)
%===============================================================================
\section{Introduction}
\label{sec:intro}

% Hook: false confidence problem
% - Existing tools give pass/fail but no indication of uncertainty
% - Under noisy conditions, this leads to false confidence
% - Real-world consequence: developers trust passing results that are meaningless

% TODO: Opening hook about false confidence in timing analysis
% Example: a developer runs dudect, sees "pass", deploys code, but the test
% was meaningless because noise swamped any potential signal

% TODO: Brief background on timing side-channels and why detection matters
% - Remote timing attacks (Timeless Timing Attacks demonstrate LAN-like precision)
% - Local attacks (SGX, containers, hyperthreading)
% - Importance in CI pipelines

% TODO: State the problem clearly
% - Binary verdicts conflate "safe" with "couldn't tell"
% - No interpretable uncertainty quantification
% - Existing tools assume iid measurements (violated in practice)

% TODO: Introduce tacet as the solution
% - Three-way decisions: Pass, Fail, Inconclusive
% - Bayesian inference provides calibrated P(leak > theta | data)
% - Block bootstrap handles autocorrelation
% - Quality gates prevent false confidence

%-------------------------------------------------------------------------------
\subsection{Contributions}
\label{sec:contributions}

% Bulleted list of contributions (usually 3-5)
% Each contribution should be concrete and verifiable

Our contributions are:

\begin{enumerate}
    \item \textbf{Calibrated verdicts}: We introduce a three-way decision framework
          (Pass/Fail/Inconclusive) that explicitly reports when evidence is
          insufficient, rather than forcing potentially incorrect verdicts.

    \item \textbf{Interpretable uncertainty}: \tacet outputs posterior probabilities
          $\prob{\text{leak} > \theta \given \text{data}}$ rather than p-values,
          providing directly interpretable uncertainty quantification.

    \item \textbf{Robustness to dependent data}: We employ block bootstrap on the
          acquisition stream to handle autocorrelation in timing measurements,
          a common violation of iid assumptions that invalidates standard
          statistical tests.

    \item \textbf{Practical tooling}: We provide a production-ready implementation
          in Rust with C/C++, Go, and JavaScript (via WASM) bindings, designed for CI integration.

    \item \textbf{Comprehensive evaluation}: We evaluate \tacet against existing
          tools (\dudect, \tlsfuzzer, \rtlf, \silent) on synthetic benchmarks
          with controlled effect sizes, autocorrelation, and noise levels,
          as well as real cryptographic libraries.
\end{enumerate}

%===============================================================================
% SECTION 2: BACKGROUND & RELATED WORK (~1.25 pages)
%===============================================================================
\section{Background and Related Work}
\label{sec:background}

% TODO: Timing side-channels 101
% - What are timing side-channels
% - Why they matter (real attacks: Lucky 13, Raccoon, MARVIN)
% - Attacker models: remote, LAN, local/shared hardware

%-------------------------------------------------------------------------------
\subsection{Timing Side-Channel Attacks}
\label{sec:background:attacks}

% Brief overview of timing attacks and their practical impact
% Mention specific attacks: Lucky 13, Raccoon, MARVIN, etc.

%-------------------------------------------------------------------------------
\subsection{Existing Detection Tools}
\label{sec:background:tools}

% Survey of existing tools with their approach and limitations

% TODO: dudect
% - Welch's t-test approach
% - Simple, fast, widely used
% - Limitations: assumes normality, no uncertainty quantification

% TODO: tlsfuzzer
% - Network-level timing analysis
% - Multiple statistical tests
% - Limitations: complex setup, autocorrelation handling limited

% TODO: RTLF (USENIX Security 2024)
% - Empirical iid bootstrap for bounded Type-1 error (10,000 iterations)
% - 9-decile comparison with per-decile thresholds
% - Binary verdict (leak/no leak based on any decile exceeding threshold)
% - Limitations:
%   * Uses iid bootstrap (does NOT handle autocorrelation)
%   * Binary verdicts only, no Inconclusive
%   * No uncertainty quantification beyond Type-I error bound
%   * R-only, operates on pre-collected measurement files
%   * No integrated measurement capability

% TODO: SILENT (arXiv 2024)
% - Non-parametric quantile-based test statistic (we adopt this)
% - m-dependent block bootstrap (DOES handle autocorrelation)
% - Bounded error rates via relevance testing
% - Limitations: binary verdicts only, acknowledges stationarity as limitation,
%   no active drift detection, R-only research tool
% - Key differentiator: Tacet outputs Inconclusive when SILENT would force a verdict

%-------------------------------------------------------------------------------
\subsection{Feature Comparison}
\label{sec:background:comparison}

% Feature comparison table (from USENIX_PAPER_PLAN.md)
% This is a key differentiator - shows tacet's unique position

\begin{table*}[t]
\centering
\caption{Feature comparison of timing side-channel detection tools.
         $\checkmark$ = supported, $\circ$ = partial support, $\times$ = not supported.
         \tacet is the only tool providing calibrated three-way verdicts with interpretable posterior probabilities.}
\label{tab:comparison}
\begin{tabular}{lccccc}
\toprule
\textbf{Feature} & \textbf{\dudect} & \textbf{\tlsfuzzer} & \textbf{\rtlf} & \textbf{\silent} & \textbf{\tacet} \\
\midrule
Calibrated verdicts\textsuperscript{1} & $\times$ & $\times$ & $\times$ & $\times$ & $\checkmark$ \\
Posterior probability\textsuperscript{2} & $\times$ & $\times$ & $\times$ & $\times$ & $\checkmark$ \\
Stationarity monitoring & $\times$ & $\times$ & $\times$ & $\times$\textsuperscript{3} & $\checkmark$ \\
Handles autocorrelation & $\times$ & $\circ$ & $\times$ & $\checkmark$\textsuperscript{4} & $\checkmark$ \\
Model-free dependence\textsuperscript{6} & $\times$ & $\times$ & $\times$ & $\times$ & $\checkmark$ \\
Non-parametric & $\times$ & $\circ$ & $\checkmark$ & $\checkmark$ & $\checkmark$ \\
Bounded false positive rate & $\circ$ & $\checkmark$ & $\checkmark$ & $\checkmark$ & $\checkmark$ \\
Coarse/discrete timers & $\checkmark$ & $\times$ & $\times$ & $\checkmark$ & $\checkmark$ \\
Effect size quantification & $\times$ & $\times$ & $\circ$\textsuperscript{5} & $\checkmark$ & $\checkmark$ \\
Adaptive sampling & $\circ$ & $\times$ & $\times$ & $\checkmark$ & $\checkmark$ \\
\midrule
Language support & C/C++/Rust & Python & R & R & C/C++/Go/JS/Rust \\
Integrated measurement & $\checkmark$ & $\checkmark$ & $\times$ & $\times$ & $\checkmark$ \\
\bottomrule
\end{tabular}
\begin{flushleft}
\small
\textsuperscript{1}Reports Inconclusive when evidence is insufficient, rather than forcing a potentially incorrect verdict.
\textsuperscript{2}Outputs interpretable $\prob{\text{leak} > \theta \given \text{data}}$ rather than binary verdict or p-value.
\textsuperscript{3}\silent acknowledges stationarity as a limitation; \tacet actively detects drift via quality gates.
\textsuperscript{4}\silent uses $m$-dependent block bootstrap but assumes a specific parametric dependence structure.
\textsuperscript{5}\rtlf reports raw decile differences without confidence intervals.
\textsuperscript{6}No parametric assumption on dependence structure; uses block bootstrap on acquisition stream without $m$-dependence model.
\end{flushleft}
\end{table*}

%===============================================================================
% SECTION 3: TACET'S APPROACH (~3.5 pages)
%===============================================================================
\section{Tacet's Approach}
\label{sec:approach}

% Problem-solution structure (from USENIX_PAPER_PLAN.md):
% | Problem with existing tools              | Tacet's solution                           |
% |------------------------------------------|-------------------------------------------|
% | Binary verdicts conflate "safe"/"unknown"| Three-way decisions with Inconclusive     |
% | No interpretable uncertainty             | Posterior probability P(leak > θ | data)  |
% | Assume iid measurements                  | Block bootstrap on acquisition stream     |
% | Can't handle drift/non-stationarity      | Condition-number monitoring, quality gates|
% | Fixed sample sizes or invalid stopping   | Adaptive sampling (likelihood-based)      |
% | No concept of resolution limits          | Explicit measurement floor θ_floor(n)     |
% | Flaky in CI under noise                  | Quality gates force Inconclusive          |

% NOTE: Design Goals merged into section intro to save space
% Key goals: calibrated verdicts, interpretable output, robust to real-world
% measurement conditions, practical CI integration

\paragraph{Threshold selection.}
The threshold $\theta$ defines the minimum effect size worth detecting. There is no
single correct value---it depends on the threat model. We provide presets based on
published attacks: \textsc{SharedHardware} (0.4\,ns, $\approx$2 cycles at 5\,GHz) for co-resident
attackers exploiting cache timing~\cite{bernstein05}; \textsc{AdjacentNetwork}
(100\,ns) for LAN attackers leveraging request multiplexing~\cite{timelesstiming2020};
\textsc{RemoteNetwork} (50\,$\mu$s) for general internet exposure. Users may
specify custom thresholds for specialized scenarios.

%-------------------------------------------------------------------------------
\subsection{Three-Way Decisions}
\label{sec:approach:decisions}

% Explain the Pass/Fail/Inconclusive framework
% - Pass: P(leak > θ | data) < 0.05
% - Fail: P(leak > θ | data) > 0.95
% - Inconclusive: evidence insufficient to reach either threshold
%
% Contrast with binary verdicts that conflate "safe" with "couldn't tell"

%-------------------------------------------------------------------------------
\subsection{Quantile-Based Test Statistic}
\label{sec:approach:quantile}

% Explain why we use quantiles (adopted from SILENT)
% - Captures both uniform shifts (branch changes) and tail effects (cache misses)
% - Non-parametric: no distributional assumptions
% - Uses 9 deciles and max_k decision functional
%
% IMPORTANT ATTRIBUTION: "We adopt the quantile-difference statistic and max-k
% decision rule from SILENT~\cite{silent2024}, embedding them in a probabilistic
% framework that enables calibrated uncertainty quantification. Our key
% contribution is not the test statistic itself, but the Bayesian inference
% framework that produces calibrated three-way verdicts with interpretable
% posterior probabilities."

%-------------------------------------------------------------------------------
\subsection{Bayesian Inference}
\label{sec:approach:inference}

% Explain the probabilistic model
% - Student-t likelihood (heavier tails than Gaussian for robustness)
% - Prior specification (MDE-based from calibration)
% - Posterior: P(leak > θ | data)
%
% Why Bayesian rather than frequentist?
% - Direct probability statements
% - Principled handling of prior information
% - Valid adaptive sampling

%-------------------------------------------------------------------------------
\subsection{Handling Autocorrelation}
\label{sec:approach:autocorrelation}

% Block bootstrap on the acquisition stream
% - Why standard bootstrap fails with dependent data (affects RTLF)
% - Block size selection: n^(1/3) heuristic (Politis-White)
% - Effective sample size calculation
%
% Note: SILENT also uses block bootstrap (m-dependent), so autocorrelation handling
% is NOT a unique differentiator. The key differences are:
% 1. Tacet doesn't assume a specific parametric m-dependence structure
% 2. Tacet actively monitors for stationarity violations (quality gates)
% 3. Tacet outputs Inconclusive when conditions degrade

%-------------------------------------------------------------------------------
\subsection{Quality Gates}
\label{sec:approach:gates}

% KEY DIFFERENTIATOR: This is what distinguishes tacet from ALL competitors
%
% Explain the quality gate mechanism:
% - Insufficient information gain (KL divergence < 0.7 nats)
% - Not learning (KL divergence collapse for 5+ batches)
% - Would take too long (projected time > 10x budget)
% - Condition drift detected (variance ratio, autocorrelation change, mean drift)
% - Condition number κ < 0.3: triggers likelihood inflation warning
%
% These gates force Inconclusive when the model can't reliably distinguish
% signal from noise - preventing false confidence.
%
% SILENT acknowledges stationarity as a limitation but does not actively detect
% violations. Tacet's quality gates fill this gap: when conditions change
% during measurement, tacet reports Inconclusive rather than producing
% a potentially misleading verdict.

%-------------------------------------------------------------------------------
\subsection{Measurement Floor}
\label{sec:approach:floor}

% Explicit θ_floor(n_eff) as a first-class concept
% - Shrinks with sqrt(n_eff)
% - Pass requires θ_eff ≤ θ_user
% - Prevents claiming "no leak" when resolution is insufficient

%-------------------------------------------------------------------------------
\subsection{Worked Example}
\label{sec:approach:example}

% ~1/3 page figure showing the same noisy dataset:
% - dudect says "Pass" (t < 10)
% - tacet says "Inconclusive" with posterior visualization
%
% This demonstrates the key value proposition: tacet doesn't give false confidence

% TODO: Create figure showing:
% 1. Timing distribution (noisy, overlapping)
% 2. dudect output: "PASS" with t-statistic
% 3. tacet output: "INCONCLUSIVE" with posterior P(leak > θ | data) ≈ 0.4

%===============================================================================
% SECTION 4: IMPLEMENTATION (~0.75 pages)
%===============================================================================
\section{Implementation}
\label{sec:implementation}

% Brief overview of the implementation
% - Rust core with unsafe blocks for timing-critical paths
% - Block bootstrap parallelized with rayon
% - C/C++ and Go bindings via FFI
% - ~X KLOC (mention size)

%-------------------------------------------------------------------------------
\subsection{Measurement Infrastructure}
\label{sec:implementation:measurement}

% Platform-specific high-resolution timing
% - x86_64: rdtsc
% - ARM64: cntvct_el0 (with adaptive batching for coarse resolution)
% - Optional PMU-based timing (kperf on macOS, perf_event on Linux)

%-------------------------------------------------------------------------------
\subsection{Adaptive Batching}
\label{sec:implementation:batching}

% On platforms with coarse timer resolution (e.g., Apple Silicon's 42ns)
% - Pilot phase to measure operation time
% - K selection (1-20) to achieve 50+ timer ticks per measurement
% - Automatically disabled on x86_64 or with PMU timers

%-------------------------------------------------------------------------------
\subsection{Bindings and CI Integration}
\label{sec:implementation:bindings}

% NOTE: Merged Language Bindings and CI Integration to save space
% Single sentence in main text: "Tacet provides C/C++, Go, and JavaScript (via WASM
% with native timers for Node.js/Deno) bindings, enabling integration into existing
% crypto library test suites and CI pipelines."
%
% Key points to cover briefly:
% - Same API surface across languages
% - Memory-safe interop
% - Configurable time budgets for CI
% - Non-flaky: Inconclusive rather than random pass/fail
% - Exit codes for CI scripts

%===============================================================================
% SECTION 5: EVALUATION (~3.5 pages)
%===============================================================================
\section{Evaluation}
\label{sec:evaluation}

% Key questions:
% 1. Does tacet provide calibrated verdicts? (FPR/FNR under varying conditions)
% 2. How does tacet compare to existing tools?
% 3. Does tacet work on real cryptographic libraries?

%-------------------------------------------------------------------------------
\subsection{Methodology}
\label{sec:evaluation:methodology}

% Experimental setup
% - Hardware configuration (CPU, memory, OS)
% - Tools compared: dudect, tlsfuzzer, RTLF, SILENT, tacet
% - Fixed sample sizes across tools for fair comparison
% - Metrics: FPR, FNR, Inconclusive rate

%-------------------------------------------------------------------------------
\subsection{Synthetic Benchmarks}
\label{sec:evaluation:synthetic}

% Three heatmaps (key evaluation figures):
%
% 1. Effect size × Tool (fixed moderate noise)
%    - Shows power comparison
%    - tacet detects leaks when present, Inconclusive when too small
%
% 2. Autocorrelation × Tool (fixed small effect)
%    - Shows robustness to dependent data
%    - Other tools break down; tacet maintains calibration
%
% 3. Noise level × Tool (null hypothesis)
%    - Shows FPR under stress
%    - Other tools show scattered false positives
%    - tacet shows Inconclusive transitioning to correct Pass

% TODO: Include heatmap figures
% Color scheme:
% - Green: Pass (correct when no leak)
% - Red: Fail/Leak detected
% - Yellow/Amber: Inconclusive
% - Gray: Error

%-------------------------------------------------------------------------------
\subsection{Results}
\label{sec:evaluation:results}

% Discussion of benchmark results
% - tacet maintains calibrated verdicts across conditions
% - Other tools: false positives under noise, false negatives under autocorrelation
% - tacet's Inconclusive rate is the tradeoff for reliability

%-------------------------------------------------------------------------------
\subsection{Crypto Library Validation}
\label{sec:evaluation:crypto}

% Real-world validation against cryptographic libraries
% - AES (constant-time implementations)
% - RSA (potential timing leaks in implementations)
% - ECC (Curve25519)
% - Post-quantum (ML-KEM, ML-DSA)
%
% Mention: incidentally rediscovered timing behavior consistent with MARVIN (CVE-2023-49092)

%-------------------------------------------------------------------------------
\subsection{Performance}
\label{sec:evaluation:performance}

% Runtime overhead comparison
% - tacet: ~100ms inference overhead per analysis
% - Slower than dudect/SILENT, faster than RTLF
% - Acceptable for CI where measurement dominates

%===============================================================================
% SECTION 6: DISCUSSION & LIMITATIONS (~0.75 pages)
%===============================================================================
\section{Discussion and Limitations}
\label{sec:discussion}

%-------------------------------------------------------------------------------
\subsection{Tradeoffs}
\label{sec:discussion:tradeoffs}

% NOTE: Combined Pipeline Complexity, Lower Power, and Overhead to save space
%
% Pipeline Complexity:
% - tacet's methodology is more complex than a t-test
% - Detailed spec required (available as artifact)
% - This complexity is the tradeoff for calibration
%
% Lower Power Under Noise:
% - tacet is more conservative
% - May require more samples to reach Pass/Fail
% - Inconclusive rate is the "price of reliability"
% - Feature, not bug: prevents false confidence
%
% Overhead:
% - ~100ms inference overhead per analysis
% - Acceptable for CI pipelines where measurement time typically dominates

%-------------------------------------------------------------------------------
\subsection{Future Work}
\label{sec:discussion:future}

% Power/EM side-channels
% - Experimental support exists
% - Different measurement infrastructure
% - Same statistical methodology applies

%===============================================================================
% SECTION 7: CONCLUSION (~0.5 pages)
%===============================================================================
\section{Conclusion}
\label{sec:conclusion}

% Summary of contributions
% - Calibrated verdicts prevent false confidence
% - Interpretable posterior probabilities
% - Robust to autocorrelation via block bootstrap
% - Practical CI integration

% Call to action
% - Available as open source (MPL v2)
% - Designed for adoption by cryptographic library maintainers

%===============================================================================
% ACKNOWLEDGMENTS (omit for submission, add for camera-ready)
%===============================================================================
% \section*{Acknowledgments}
% Omitted for anonymous submission.

%===============================================================================
% REQUIRED APPENDICES
%===============================================================================
\cleardoublepage
\appendix

%-------------------------------------------------------------------------------
\section*{Ethical Considerations}
\label{app:ethics}

% Stakeholder-based analysis following Menlo Report principles

\textbf{Stakeholders.}
We identify three primary stakeholders:
(1) cryptographic library maintainers, who benefit from improved vulnerability detection;
(2) end users of cryptographic software, who benefit from more secure implementations; and
(3) potential attackers, who could misuse timing analysis techniques.

\textbf{Impacts.}
Following the Menlo Report principles:

\begin{itemize}
    \item \emph{Beneficence}: \tacet enables proactive detection of timing
          side-channels before deployment, reducing the window of vulnerability
          for end users.

    \item \emph{Respect for Persons}: Our work does not involve human subjects
          or personal data.

    \item \emph{Justice}: By releasing \tacet as open-source software (MPL v2),
          we ensure equal access to improved security testing across
          organizations of all sizes.

    \item \emph{Respect for Law and Public Interest}: Timing side-channel
          detection is an established defensive practice; we are not
          introducing novel attack capabilities.
\end{itemize}

\textbf{Potential Harms.}
The statistical techniques underlying \tacet could theoretically help attackers
confirm timing leaks. However, simpler tools (\eg \dudect) already provide
this capability; \tacet's contribution is primarily in reducing false positives
and providing calibrated uncertainty---properties more valuable to defenders
than attackers.

\textbf{Decision.}
We concluded that publication serves the public interest: \tacet's practical
CI integration lowers the barrier for maintainers to adopt rigorous timing
analysis, catching vulnerabilities before they affect users. The marginal
benefit to attackers is minimal compared to existing tools.

%-------------------------------------------------------------------------------
\cleardoublepage
\section*{Open Science}
\label{app:open-science}

All artifacts necessary to evaluate our contributions are publicly available:

\begin{itemize}
    \item \textbf{Source code}: \url{https://anonymous.4open.science/r/tacet-75EB/}
          (MPL v2 license)

    \item \textbf{Specification}: Included in repository as
          \texttt{website/src/content/reference/specification.md}

    \item \textbf{Benchmark scripts}: \texttt{crates/tacet-bench/} directory,
          with instructions in \texttt{README.md}

    \item \textbf{Benchmark results}: Raw data available in artifact
\end{itemize}

The repository includes instructions for reproducing all experiments reported
in Section~\ref{sec:evaluation}. Hardware configuration details are provided
in Section~\ref{sec:evaluation:methodology}.

%===============================================================================
% OPTIONAL APPENDIX: STATISTICAL METHODOLOGY
%===============================================================================
\cleardoublepage
% Methodological Appendix for Tacet USENIX Security 2026 Paper
% To be included via % Methodological Appendix for Tacet USENIX Security 2026 Paper
% To be included via % Methodological Appendix for Tacet USENIX Security 2026 Paper
% To be included via \input{appendix_methodology} in the main document

\section{Statistical Methodology Details}
\label{app:methodology}

This appendix provides additional detail on tacet's statistical methodology for reviewers familiar with Bayesian inference. We cover the probabilistic model, inference procedure, autocorrelation handling, quality gates, and measurement floor derivation. For complete algorithmic specifications, see the specification document in the artifact repository.

\subsection{Probabilistic Model}
\label{app:model}

\paragraph{Test Statistic.}
Rather than comparing means (as in t-tests), we compare timing distributions via their deciles. Let $F$ and $R$ denote timing samples from fixed (e.g., all-zeros) and random input classes, respectively. The test statistic is the vector of quantile differences:
\begin{equation}
\Delta = \hat{q}(F) - \hat{q}(R) \in \mathbb{R}^9
\end{equation}
where $\hat{q}_p$ denotes the empirical $p$-th quantile for $p \in \{0.1, 0.2, \ldots, 0.9\}$. This captures both uniform shifts (all quantiles affected equally, indicating branch changes) and tail effects (upper quantiles affected more, indicating cache-based leaks).

The decision functional is $m(\delta) = \max_k |\delta_k|$---the maximum absolute quantile difference. A timing leak exceeding threshold $\theta$ is declared when $m(\delta) > \theta$.

\paragraph{Likelihood.}
Quantile estimators are asymptotically normal, but real timing data can deviate from this approximation when covariance is underestimated. To prevent pathological certainty under covariance misspecification, we use a robust Student-$t$ likelihood via a scale-mixture representation:
\begin{align}
\Delta \mid \delta, \kappa &\sim \mathcal{N}\!\left(\delta, \frac{\Sigma_n}{\kappa}\right) \\
\kappa &\sim \text{Gamma}\!\left(\frac{\nu_\ell}{2}, \frac{\nu_\ell}{2}\right)
\end{align}
with $\nu_\ell = 8$. Marginally, this yields $\Delta \mid \delta \sim t_{\nu_\ell}(\delta, \Sigma_n)$.

When the estimated covariance $\Sigma_n$ is too small (e.g., due to unmodeled autocorrelation), $\kappa$ is pulled below 1, automatically inflating uncertainty. This is more robust than a Gaussian likelihood, which would produce overconfident posteriors under covariance misestimation.

\paragraph{Prior.}
We place a multivariate Student-$t$ prior on the latent effect profile $\delta$:
\begin{equation}
\delta \sim t_\nu(0, \sigma^2 R), \quad \nu = 4
\end{equation}
where $R = \text{Corr}(\Sigma_{\text{rate}})$ is the prior correlation matrix derived from the covariance rate estimated during calibration. The scale $\sigma$ is calibrated so that the prior probability of exceeding the effective threshold $\theta_{\text{eff}}$ equals a target $\pi_0 = 0.62$:
\begin{equation}
P\!\left(\max_k |\delta_k| > \theta_{\text{eff}} \mid \delta \sim t_\nu(0, \sigma^2 R)\right) = \pi_0
\end{equation}
This is solved by deterministic 1D root-finding using Monte Carlo integration (50,000 samples). The resulting prior expresses genuine uncertainty: it does not strongly favor either ``leak'' or ``no leak'' before seeing data.

We implement inference using the equivalent scale-mixture representation:
\begin{align}
\lambda &\sim \text{Gamma}\!\left(\frac{\nu}{2}, \frac{\nu}{2}\right) \\
\delta \mid \lambda &\sim \mathcal{N}\!\left(0, \frac{\sigma^2}{\lambda} R\right)
\end{align}
which enables closed-form Gibbs conditionals.

\subsection{Inference Procedure}
\label{app:inference}

\paragraph{Gibbs Sampler.}
We approximate the posterior $p(\delta \mid \Delta)$ using a deterministic Gibbs sampler over $(\delta, \lambda, \kappa)$. Each iteration samples:
\begin{enumerate}
\item $\delta^{(t)} \sim p(\delta \mid \lambda^{(t-1)}, \kappa^{(t-1)}, \Delta)$ --- a multivariate normal
\item $\lambda^{(t)} \sim p(\lambda \mid \delta^{(t)})$ --- a gamma distribution
\item $\kappa^{(t)} \sim p(\kappa \mid \delta^{(t)}, \Delta)$ --- a gamma distribution
\end{enumerate}

All conditionals are conjugate and admit closed-form sampling. The $\delta$ conditional requires computing a posterior precision matrix $Q = \kappa \Sigma_n^{-1} + (\lambda/\sigma^2) R^{-1}$ and sampling from $\mathcal{N}(Q^{-1}(\kappa \Sigma_n^{-1} \Delta), Q^{-1})$. All matrix operations use Cholesky decomposition for numerical stability---we never form explicit inverses.

\paragraph{Convergence and Cost.}
We run 256 total iterations with 64 burn-in, retaining 192 samples. This is sufficient because: (1) the model is low-dimensional ($d=9$) with conjugate structure, ensuring fast mixing; (2) initialization at $\lambda^{(0)} = \kappa^{(0)} = 1$ (the prior means) provides a good starting point; and (3) the decision (leak probability) depends only on whether $m(\delta^{(s)}) > \theta_{\text{eff}}$, which is robust to modest posterior approximation error.

We track $\lambda$ mixing via its coefficient of variation (CV) and effective sample size (ESS). If $\lambda$ CV $< 0.1$ or $\lambda$ ESS $< 20$, we flag a potential convergence issue in diagnostics, though this rarely occurs in practice.

Computational cost per posterior update is $O(d^3 + d^2 \cdot N_{\text{gibbs}})$, dominated by the Cholesky decompositions. With $d=9$ and $N_{\text{gibbs}}=256$, this takes $<1$ms on modern hardware---negligible compared to measurement time.

\paragraph{Leak Probability.}
The leak probability is estimated as the Monte Carlo fraction of posterior samples exceeding the threshold:
\begin{equation}
\widehat{P}(\text{leak} > \theta_{\text{eff}} \mid \Delta) = \frac{1}{N_{\text{keep}}} \sum_{s=1}^{N_{\text{keep}}} \mathbf{1}\!\left[m(\delta^{(s)}) > \theta_{\text{eff}}\right]
\end{equation}

\subsection{Block Bootstrap for Autocorrelation}
\label{app:bootstrap}

\paragraph{Why Standard Bootstrap Fails.}
Timing measurements exhibit autocorrelation: nearby samples are more similar than distant ones due to cache state, CPU frequency scaling, and OS scheduling. Standard bootstrap assumes i.i.d.\ samples, systematically underestimating variance and producing overconfident posteriors.

For example, if consecutive measurements have autocorrelation $\rho = 0.3$, the effective sample size may be 30--50\% lower than the nominal count. Ignoring this inflates the false positive rate---the main failure mode we aim to prevent.

\paragraph{Stream-Based Block Bootstrap.}
We use block bootstrap on the \emph{acquisition stream}---the interleaved sequence of (class, timing) pairs in measurement order---rather than per-class samples. This preserves the true temporal dependence structure.

The algorithm resamples contiguous blocks of length $\hat{b}$ with replacement, reconstructs a bootstrap stream of the same length, splits by class, and computes quantile differences. Over $B = 2000$ iterations, we estimate the covariance matrix $\Sigma_{\text{cal}}$ of $\Delta$.

\paragraph{Block Length Selection.}
We use the Politis-White automatic block length selection algorithm with a class-conditional modification. Standard Politis-White estimates autocorrelation from raw lag-$k$ pairs, but with interleaved sampling, a lag-$k$ pair may cross class boundaries. Instead, we compute autocorrelation using only same-class pairs at each acquisition-stream lag, then take the conservative maximum across classes.

The optimal block length scales as $\hat{b} \propto T^{1/3}$ where $T$ is the stream length. We clamp $\hat{b} \in [10, \min(3\sqrt{T}, T/3)]$ to prevent degenerate blocks. In high-autocorrelation regimes (detected when $\rho_{\text{max}}(10) > 0.3$), we inflate the block length by 50\% for additional conservatism.

\paragraph{Effective Sample Size via IACT.}
We compute the integrated autocorrelation time (IACT) $\hat{\tau}$ using the Geyer initial monotone sequence (IMS) algorithm. This estimates how many nominal samples equal one independent sample:
\begin{equation}
n_{\text{eff}} = \frac{n}{\hat{\tau}}
\end{equation}

For timing analysis, IACT is computed on indicator series $\mathbf{1}[y \leq q_p]$ for each quantile $p$, taking the maximum across quantiles and classes. This captures dependence in the quantile estimators, not just raw timings.

\paragraph{Covariance Scaling.}
The key insight enabling adaptive sampling is that quantile estimator covariance scales as $1/n$. We estimate a ``covariance rate'' during calibration:
\begin{equation}
\Sigma_{\text{rate}} = \hat{\Sigma}_{\text{cal}} \cdot n_{\text{cal}}
\end{equation}
During the adaptive loop with $n$ samples and current IACT $\hat{\tau}$:
\begin{equation}
\Sigma_n = \Sigma_{\text{rate}} \cdot \frac{\hat{\tau}}{\hat{\tau}_{\text{cal}}} \cdot \frac{1}{n}
\end{equation}
This allows cheap posterior updates without re-bootstrapping at each batch.

\subsection{Quality Gates}
\label{app:gates}

Quality gates detect when data quality is insufficient for a confident verdict, forcing an \textsc{Inconclusive} outcome rather than an unreliable Pass or Fail. This is the key mechanism ensuring calibrated verdicts.

\paragraph{Gate 1: Insufficient Information Gain.}
We compute the KL divergence between Gaussian approximations of the prior and posterior:
\begin{multline}
\text{KL}(p_{\text{post}} \| p_{\text{prior}}) = \frac{1}{2}\Big(
\text{tr}(\Lambda_0^{-1}\Lambda_{\text{post}}) \\
+ \mu_{\text{post}}^\top \Lambda_0^{-1} \mu_{\text{post}}
- d + \ln\frac{|\Lambda_0|}{|\Lambda_{\text{post}}|}
\Big)
\end{multline}
where $\Lambda_0 = 2\sigma^2 R$ is the marginal prior covariance for $\nu=4$, and $(\mu_{\text{post}}, \Lambda_{\text{post}})$ are estimated from Gibbs samples.

\textbf{DataTooNoisy}: If KL $< 0.7$ nats after calibration, the posterior has barely moved from the prior---the data is too noisy to distinguish signal from noise.

\textbf{NotLearning}: If the sum of inter-batch KL divergences falls below 0.001 for 5+ consecutive batches, the posterior has stopped updating. Additional samples will not help.

\paragraph{Gate 2: Would Take Too Long.}
We extrapolate time to decision based on current convergence rate. If projected time exceeds 10$\times$ the configured budget, we stop early with \textsc{Inconclusive} rather than waste resources on a test that cannot conclude.

\paragraph{Gate 3: Resource Budget Exceeded.}
If elapsed time exceeds the configured time budget, or total samples exceed the configured maximum, we stop with \textsc{Inconclusive} and report the current posterior probability.

\paragraph{Gate 4: Condition Drift Detected.}
The covariance rate $\Sigma_{\text{rate}}$ is estimated during calibration. If measurement conditions change during the adaptive loop, this estimate becomes invalid. We detect drift by comparing:
\begin{itemize}
\item Variance ratio: $\sigma_{\text{post}}^2 / \sigma_{\text{cal}}^2$ outside $[0.5, 2.0]$
\item Autocorrelation change: $|\rho_{\text{post}}(1) - \rho_{\text{cal}}(1)| > 0.3$
\item Mean drift: $|\mu_{\text{post}} - \mu_{\text{cal}}| / \sigma_{\text{cal}} > 3.0$
\end{itemize}
Any trigger yields \textsc{Inconclusive} with reason \textsc{ConditionsChanged}.

\paragraph{Likelihood Inflation Warning.}
When the posterior mean of $\kappa$ falls below 0.3, the likelihood covariance was effectively scaled up by $>3\times$ for robustness. This indicates potential covariance misestimation and is reported in diagnostics, though it does not block verdicts.

\subsection{Measurement Floor Derivation}
\label{app:floor}

A critical design element is distinguishing what the user \emph{wants} to detect ($\theta_{\text{user}}$) from what the measurement \emph{can} detect ($\theta_{\text{floor}}$).

\paragraph{Floor-Rate Constant.}
During calibration, we compute a floor-rate constant $c_{\text{floor}}$ representing the 95th percentile of the null distribution of $m(\delta)$:
\begin{equation}
c_{\text{floor}} = q_{0.95}\!\left(\max_k |Z_{0,k}|\right), \quad Z_0 \sim \mathcal{N}(0, \Sigma_{\text{rate}})
\end{equation}
estimated via Monte Carlo (50,000 samples). This characterizes the ``noise floor'' of the measurement system.

\paragraph{Dynamic Measurement Floor.}
During the adaptive loop, the statistical floor shrinks with effective sample size:
\begin{equation}
\theta_{\text{floor,stat}}(n) = \frac{c_{\text{floor}}}{\sqrt{n_{\text{eff}}}}
\end{equation}

The timer also imposes a floor based on quantization:
\begin{equation}
\theta_{\text{tick}} = \frac{\text{1 tick (ns)}}{K}
\end{equation}
where $K$ is the batch size (if batching is used). The combined floor is:
\begin{equation}
\theta_{\text{floor}}(n) = \max\!\left(\theta_{\text{floor,stat}}(n), \theta_{\text{tick}}\right)
\end{equation}

\paragraph{Effective Threshold.}
The effective threshold used for inference is:
\begin{equation}
\theta_{\text{eff}} = \max(\theta_{\text{user}}, \theta_{\text{floor}})
\end{equation}

When $\theta_{\text{eff}} > \theta_{\text{user}}$:
\begin{itemize}
\item \textbf{Fail propagates}: Detecting $m(\delta) > \theta_{\text{eff}}$ implies $m(\delta) > \theta_{\text{user}}$, since $\theta_{\text{eff}} \geq \theta_{\text{user}}$.
\item \textbf{Pass does not propagate}: ``No effect above $\theta_{\text{eff}}$'' is compatible with effects in $(\theta_{\text{user}}, \theta_{\text{eff}}]$. We return \textsc{Inconclusive} with reason \textsc{ThresholdElevated}.
\end{itemize}

As sampling continues, $\theta_{\text{floor}}(n)$ decreases. If it drops to $\theta_{\text{user}}$ or below, Pass becomes possible (subject to the posterior). This is why adaptive sampling can eventually reach Pass even when early calibration shows an elevated floor.

\vspace{1em}
\noindent\textit{For complete algorithmic details including pseudocode, numerical stability procedures, and platform-specific considerations, see the specification document in the artifact repository.}
 in the main document

\section{Statistical Methodology Details}
\label{app:methodology}

This appendix provides additional detail on tacet's statistical methodology for reviewers familiar with Bayesian inference. We cover the probabilistic model, inference procedure, autocorrelation handling, quality gates, and measurement floor derivation. For complete algorithmic specifications, see the specification document in the artifact repository.

\subsection{Probabilistic Model}
\label{app:model}

\paragraph{Test Statistic.}
Rather than comparing means (as in t-tests), we compare timing distributions via their deciles. Let $F$ and $R$ denote timing samples from fixed (e.g., all-zeros) and random input classes, respectively. The test statistic is the vector of quantile differences:
\begin{equation}
\Delta = \hat{q}(F) - \hat{q}(R) \in \mathbb{R}^9
\end{equation}
where $\hat{q}_p$ denotes the empirical $p$-th quantile for $p \in \{0.1, 0.2, \ldots, 0.9\}$. This captures both uniform shifts (all quantiles affected equally, indicating branch changes) and tail effects (upper quantiles affected more, indicating cache-based leaks).

The decision functional is $m(\delta) = \max_k |\delta_k|$---the maximum absolute quantile difference. A timing leak exceeding threshold $\theta$ is declared when $m(\delta) > \theta$.

\paragraph{Likelihood.}
Quantile estimators are asymptotically normal, but real timing data can deviate from this approximation when covariance is underestimated. To prevent pathological certainty under covariance misspecification, we use a robust Student-$t$ likelihood via a scale-mixture representation:
\begin{align}
\Delta \mid \delta, \kappa &\sim \mathcal{N}\!\left(\delta, \frac{\Sigma_n}{\kappa}\right) \\
\kappa &\sim \text{Gamma}\!\left(\frac{\nu_\ell}{2}, \frac{\nu_\ell}{2}\right)
\end{align}
with $\nu_\ell = 8$. Marginally, this yields $\Delta \mid \delta \sim t_{\nu_\ell}(\delta, \Sigma_n)$.

When the estimated covariance $\Sigma_n$ is too small (e.g., due to unmodeled autocorrelation), $\kappa$ is pulled below 1, automatically inflating uncertainty. This is more robust than a Gaussian likelihood, which would produce overconfident posteriors under covariance misestimation.

\paragraph{Prior.}
We place a multivariate Student-$t$ prior on the latent effect profile $\delta$:
\begin{equation}
\delta \sim t_\nu(0, \sigma^2 R), \quad \nu = 4
\end{equation}
where $R = \text{Corr}(\Sigma_{\text{rate}})$ is the prior correlation matrix derived from the covariance rate estimated during calibration. The scale $\sigma$ is calibrated so that the prior probability of exceeding the effective threshold $\theta_{\text{eff}}$ equals a target $\pi_0 = 0.62$:
\begin{equation}
P\!\left(\max_k |\delta_k| > \theta_{\text{eff}} \mid \delta \sim t_\nu(0, \sigma^2 R)\right) = \pi_0
\end{equation}
This is solved by deterministic 1D root-finding using Monte Carlo integration (50,000 samples). The resulting prior expresses genuine uncertainty: it does not strongly favor either ``leak'' or ``no leak'' before seeing data.

We implement inference using the equivalent scale-mixture representation:
\begin{align}
\lambda &\sim \text{Gamma}\!\left(\frac{\nu}{2}, \frac{\nu}{2}\right) \\
\delta \mid \lambda &\sim \mathcal{N}\!\left(0, \frac{\sigma^2}{\lambda} R\right)
\end{align}
which enables closed-form Gibbs conditionals.

\subsection{Inference Procedure}
\label{app:inference}

\paragraph{Gibbs Sampler.}
We approximate the posterior $p(\delta \mid \Delta)$ using a deterministic Gibbs sampler over $(\delta, \lambda, \kappa)$. Each iteration samples:
\begin{enumerate}
\item $\delta^{(t)} \sim p(\delta \mid \lambda^{(t-1)}, \kappa^{(t-1)}, \Delta)$ --- a multivariate normal
\item $\lambda^{(t)} \sim p(\lambda \mid \delta^{(t)})$ --- a gamma distribution
\item $\kappa^{(t)} \sim p(\kappa \mid \delta^{(t)}, \Delta)$ --- a gamma distribution
\end{enumerate}

All conditionals are conjugate and admit closed-form sampling. The $\delta$ conditional requires computing a posterior precision matrix $Q = \kappa \Sigma_n^{-1} + (\lambda/\sigma^2) R^{-1}$ and sampling from $\mathcal{N}(Q^{-1}(\kappa \Sigma_n^{-1} \Delta), Q^{-1})$. All matrix operations use Cholesky decomposition for numerical stability---we never form explicit inverses.

\paragraph{Convergence and Cost.}
We run 256 total iterations with 64 burn-in, retaining 192 samples. This is sufficient because: (1) the model is low-dimensional ($d=9$) with conjugate structure, ensuring fast mixing; (2) initialization at $\lambda^{(0)} = \kappa^{(0)} = 1$ (the prior means) provides a good starting point; and (3) the decision (leak probability) depends only on whether $m(\delta^{(s)}) > \theta_{\text{eff}}$, which is robust to modest posterior approximation error.

We track $\lambda$ mixing via its coefficient of variation (CV) and effective sample size (ESS). If $\lambda$ CV $< 0.1$ or $\lambda$ ESS $< 20$, we flag a potential convergence issue in diagnostics, though this rarely occurs in practice.

Computational cost per posterior update is $O(d^3 + d^2 \cdot N_{\text{gibbs}})$, dominated by the Cholesky decompositions. With $d=9$ and $N_{\text{gibbs}}=256$, this takes $<1$ms on modern hardware---negligible compared to measurement time.

\paragraph{Leak Probability.}
The leak probability is estimated as the Monte Carlo fraction of posterior samples exceeding the threshold:
\begin{equation}
\widehat{P}(\text{leak} > \theta_{\text{eff}} \mid \Delta) = \frac{1}{N_{\text{keep}}} \sum_{s=1}^{N_{\text{keep}}} \mathbf{1}\!\left[m(\delta^{(s)}) > \theta_{\text{eff}}\right]
\end{equation}

\subsection{Block Bootstrap for Autocorrelation}
\label{app:bootstrap}

\paragraph{Why Standard Bootstrap Fails.}
Timing measurements exhibit autocorrelation: nearby samples are more similar than distant ones due to cache state, CPU frequency scaling, and OS scheduling. Standard bootstrap assumes i.i.d.\ samples, systematically underestimating variance and producing overconfident posteriors.

For example, if consecutive measurements have autocorrelation $\rho = 0.3$, the effective sample size may be 30--50\% lower than the nominal count. Ignoring this inflates the false positive rate---the main failure mode we aim to prevent.

\paragraph{Stream-Based Block Bootstrap.}
We use block bootstrap on the \emph{acquisition stream}---the interleaved sequence of (class, timing) pairs in measurement order---rather than per-class samples. This preserves the true temporal dependence structure.

The algorithm resamples contiguous blocks of length $\hat{b}$ with replacement, reconstructs a bootstrap stream of the same length, splits by class, and computes quantile differences. Over $B = 2000$ iterations, we estimate the covariance matrix $\Sigma_{\text{cal}}$ of $\Delta$.

\paragraph{Block Length Selection.}
We use the Politis-White automatic block length selection algorithm with a class-conditional modification. Standard Politis-White estimates autocorrelation from raw lag-$k$ pairs, but with interleaved sampling, a lag-$k$ pair may cross class boundaries. Instead, we compute autocorrelation using only same-class pairs at each acquisition-stream lag, then take the conservative maximum across classes.

The optimal block length scales as $\hat{b} \propto T^{1/3}$ where $T$ is the stream length. We clamp $\hat{b} \in [10, \min(3\sqrt{T}, T/3)]$ to prevent degenerate blocks. In high-autocorrelation regimes (detected when $\rho_{\text{max}}(10) > 0.3$), we inflate the block length by 50\% for additional conservatism.

\paragraph{Effective Sample Size via IACT.}
We compute the integrated autocorrelation time (IACT) $\hat{\tau}$ using the Geyer initial monotone sequence (IMS) algorithm. This estimates how many nominal samples equal one independent sample:
\begin{equation}
n_{\text{eff}} = \frac{n}{\hat{\tau}}
\end{equation}

For timing analysis, IACT is computed on indicator series $\mathbf{1}[y \leq q_p]$ for each quantile $p$, taking the maximum across quantiles and classes. This captures dependence in the quantile estimators, not just raw timings.

\paragraph{Covariance Scaling.}
The key insight enabling adaptive sampling is that quantile estimator covariance scales as $1/n$. We estimate a ``covariance rate'' during calibration:
\begin{equation}
\Sigma_{\text{rate}} = \hat{\Sigma}_{\text{cal}} \cdot n_{\text{cal}}
\end{equation}
During the adaptive loop with $n$ samples and current IACT $\hat{\tau}$:
\begin{equation}
\Sigma_n = \Sigma_{\text{rate}} \cdot \frac{\hat{\tau}}{\hat{\tau}_{\text{cal}}} \cdot \frac{1}{n}
\end{equation}
This allows cheap posterior updates without re-bootstrapping at each batch.

\subsection{Quality Gates}
\label{app:gates}

Quality gates detect when data quality is insufficient for a confident verdict, forcing an \textsc{Inconclusive} outcome rather than an unreliable Pass or Fail. This is the key mechanism ensuring calibrated verdicts.

\paragraph{Gate 1: Insufficient Information Gain.}
We compute the KL divergence between Gaussian approximations of the prior and posterior:
\begin{multline}
\text{KL}(p_{\text{post}} \| p_{\text{prior}}) = \frac{1}{2}\Big(
\text{tr}(\Lambda_0^{-1}\Lambda_{\text{post}}) \\
+ \mu_{\text{post}}^\top \Lambda_0^{-1} \mu_{\text{post}}
- d + \ln\frac{|\Lambda_0|}{|\Lambda_{\text{post}}|}
\Big)
\end{multline}
where $\Lambda_0 = 2\sigma^2 R$ is the marginal prior covariance for $\nu=4$, and $(\mu_{\text{post}}, \Lambda_{\text{post}})$ are estimated from Gibbs samples.

\textbf{DataTooNoisy}: If KL $< 0.7$ nats after calibration, the posterior has barely moved from the prior---the data is too noisy to distinguish signal from noise.

\textbf{NotLearning}: If the sum of inter-batch KL divergences falls below 0.001 for 5+ consecutive batches, the posterior has stopped updating. Additional samples will not help.

\paragraph{Gate 2: Would Take Too Long.}
We extrapolate time to decision based on current convergence rate. If projected time exceeds 10$\times$ the configured budget, we stop early with \textsc{Inconclusive} rather than waste resources on a test that cannot conclude.

\paragraph{Gate 3: Resource Budget Exceeded.}
If elapsed time exceeds the configured time budget, or total samples exceed the configured maximum, we stop with \textsc{Inconclusive} and report the current posterior probability.

\paragraph{Gate 4: Condition Drift Detected.}
The covariance rate $\Sigma_{\text{rate}}$ is estimated during calibration. If measurement conditions change during the adaptive loop, this estimate becomes invalid. We detect drift by comparing:
\begin{itemize}
\item Variance ratio: $\sigma_{\text{post}}^2 / \sigma_{\text{cal}}^2$ outside $[0.5, 2.0]$
\item Autocorrelation change: $|\rho_{\text{post}}(1) - \rho_{\text{cal}}(1)| > 0.3$
\item Mean drift: $|\mu_{\text{post}} - \mu_{\text{cal}}| / \sigma_{\text{cal}} > 3.0$
\end{itemize}
Any trigger yields \textsc{Inconclusive} with reason \textsc{ConditionsChanged}.

\paragraph{Likelihood Inflation Warning.}
When the posterior mean of $\kappa$ falls below 0.3, the likelihood covariance was effectively scaled up by $>3\times$ for robustness. This indicates potential covariance misestimation and is reported in diagnostics, though it does not block verdicts.

\subsection{Measurement Floor Derivation}
\label{app:floor}

A critical design element is distinguishing what the user \emph{wants} to detect ($\theta_{\text{user}}$) from what the measurement \emph{can} detect ($\theta_{\text{floor}}$).

\paragraph{Floor-Rate Constant.}
During calibration, we compute a floor-rate constant $c_{\text{floor}}$ representing the 95th percentile of the null distribution of $m(\delta)$:
\begin{equation}
c_{\text{floor}} = q_{0.95}\!\left(\max_k |Z_{0,k}|\right), \quad Z_0 \sim \mathcal{N}(0, \Sigma_{\text{rate}})
\end{equation}
estimated via Monte Carlo (50,000 samples). This characterizes the ``noise floor'' of the measurement system.

\paragraph{Dynamic Measurement Floor.}
During the adaptive loop, the statistical floor shrinks with effective sample size:
\begin{equation}
\theta_{\text{floor,stat}}(n) = \frac{c_{\text{floor}}}{\sqrt{n_{\text{eff}}}}
\end{equation}

The timer also imposes a floor based on quantization:
\begin{equation}
\theta_{\text{tick}} = \frac{\text{1 tick (ns)}}{K}
\end{equation}
where $K$ is the batch size (if batching is used). The combined floor is:
\begin{equation}
\theta_{\text{floor}}(n) = \max\!\left(\theta_{\text{floor,stat}}(n), \theta_{\text{tick}}\right)
\end{equation}

\paragraph{Effective Threshold.}
The effective threshold used for inference is:
\begin{equation}
\theta_{\text{eff}} = \max(\theta_{\text{user}}, \theta_{\text{floor}})
\end{equation}

When $\theta_{\text{eff}} > \theta_{\text{user}}$:
\begin{itemize}
\item \textbf{Fail propagates}: Detecting $m(\delta) > \theta_{\text{eff}}$ implies $m(\delta) > \theta_{\text{user}}$, since $\theta_{\text{eff}} \geq \theta_{\text{user}}$.
\item \textbf{Pass does not propagate}: ``No effect above $\theta_{\text{eff}}$'' is compatible with effects in $(\theta_{\text{user}}, \theta_{\text{eff}}]$. We return \textsc{Inconclusive} with reason \textsc{ThresholdElevated}.
\end{itemize}

As sampling continues, $\theta_{\text{floor}}(n)$ decreases. If it drops to $\theta_{\text{user}}$ or below, Pass becomes possible (subject to the posterior). This is why adaptive sampling can eventually reach Pass even when early calibration shows an elevated floor.

\vspace{1em}
\noindent\textit{For complete algorithmic details including pseudocode, numerical stability procedures, and platform-specific considerations, see the specification document in the artifact repository.}
 in the main document

\section{Statistical Methodology Details}
\label{app:methodology}

This appendix provides additional detail on tacet's statistical methodology for reviewers familiar with Bayesian inference. We cover the probabilistic model, inference procedure, autocorrelation handling, quality gates, and measurement floor derivation. For complete algorithmic specifications, see the specification document in the artifact repository.

\subsection{Probabilistic Model}
\label{app:model}

\paragraph{Test Statistic.}
Rather than comparing means (as in t-tests), we compare timing distributions via their deciles. Let $F$ and $R$ denote timing samples from fixed (e.g., all-zeros) and random input classes, respectively. The test statistic is the vector of quantile differences:
\begin{equation}
\Delta = \hat{q}(F) - \hat{q}(R) \in \mathbb{R}^9
\end{equation}
where $\hat{q}_p$ denotes the empirical $p$-th quantile for $p \in \{0.1, 0.2, \ldots, 0.9\}$. This captures both uniform shifts (all quantiles affected equally, indicating branch changes) and tail effects (upper quantiles affected more, indicating cache-based leaks).

The decision functional is $m(\delta) = \max_k |\delta_k|$---the maximum absolute quantile difference. A timing leak exceeding threshold $\theta$ is declared when $m(\delta) > \theta$.

\paragraph{Likelihood.}
Quantile estimators are asymptotically normal, but real timing data can deviate from this approximation when covariance is underestimated. To prevent pathological certainty under covariance misspecification, we use a robust Student-$t$ likelihood via a scale-mixture representation:
\begin{align}
\Delta \mid \delta, \kappa &\sim \mathcal{N}\!\left(\delta, \frac{\Sigma_n}{\kappa}\right) \\
\kappa &\sim \text{Gamma}\!\left(\frac{\nu_\ell}{2}, \frac{\nu_\ell}{2}\right)
\end{align}
with $\nu_\ell = 8$. Marginally, this yields $\Delta \mid \delta \sim t_{\nu_\ell}(\delta, \Sigma_n)$.

When the estimated covariance $\Sigma_n$ is too small (e.g., due to unmodeled autocorrelation), $\kappa$ is pulled below 1, automatically inflating uncertainty. This is more robust than a Gaussian likelihood, which would produce overconfident posteriors under covariance misestimation.

\paragraph{Prior.}
We place a multivariate Student-$t$ prior on the latent effect profile $\delta$:
\begin{equation}
\delta \sim t_\nu(0, \sigma^2 R), \quad \nu = 4
\end{equation}
where $R = \text{Corr}(\Sigma_{\text{rate}})$ is the prior correlation matrix derived from the covariance rate estimated during calibration. The scale $\sigma$ is calibrated so that the prior probability of exceeding the effective threshold $\theta_{\text{eff}}$ equals a target $\pi_0 = 0.62$:
\begin{equation}
P\!\left(\max_k |\delta_k| > \theta_{\text{eff}} \mid \delta \sim t_\nu(0, \sigma^2 R)\right) = \pi_0
\end{equation}
This is solved by deterministic 1D root-finding using Monte Carlo integration (50,000 samples). The resulting prior expresses genuine uncertainty: it does not strongly favor either ``leak'' or ``no leak'' before seeing data.

We implement inference using the equivalent scale-mixture representation:
\begin{align}
\lambda &\sim \text{Gamma}\!\left(\frac{\nu}{2}, \frac{\nu}{2}\right) \\
\delta \mid \lambda &\sim \mathcal{N}\!\left(0, \frac{\sigma^2}{\lambda} R\right)
\end{align}
which enables closed-form Gibbs conditionals.

\subsection{Inference Procedure}
\label{app:inference}

\paragraph{Gibbs Sampler.}
We approximate the posterior $p(\delta \mid \Delta)$ using a deterministic Gibbs sampler over $(\delta, \lambda, \kappa)$. Each iteration samples:
\begin{enumerate}
\item $\delta^{(t)} \sim p(\delta \mid \lambda^{(t-1)}, \kappa^{(t-1)}, \Delta)$ --- a multivariate normal
\item $\lambda^{(t)} \sim p(\lambda \mid \delta^{(t)})$ --- a gamma distribution
\item $\kappa^{(t)} \sim p(\kappa \mid \delta^{(t)}, \Delta)$ --- a gamma distribution
\end{enumerate}

All conditionals are conjugate and admit closed-form sampling. The $\delta$ conditional requires computing a posterior precision matrix $Q = \kappa \Sigma_n^{-1} + (\lambda/\sigma^2) R^{-1}$ and sampling from $\mathcal{N}(Q^{-1}(\kappa \Sigma_n^{-1} \Delta), Q^{-1})$. All matrix operations use Cholesky decomposition for numerical stability---we never form explicit inverses.

\paragraph{Convergence and Cost.}
We run 256 total iterations with 64 burn-in, retaining 192 samples. This is sufficient because: (1) the model is low-dimensional ($d=9$) with conjugate structure, ensuring fast mixing; (2) initialization at $\lambda^{(0)} = \kappa^{(0)} = 1$ (the prior means) provides a good starting point; and (3) the decision (leak probability) depends only on whether $m(\delta^{(s)}) > \theta_{\text{eff}}$, which is robust to modest posterior approximation error.

We track $\lambda$ mixing via its coefficient of variation (CV) and effective sample size (ESS). If $\lambda$ CV $< 0.1$ or $\lambda$ ESS $< 20$, we flag a potential convergence issue in diagnostics, though this rarely occurs in practice.

Computational cost per posterior update is $O(d^3 + d^2 \cdot N_{\text{gibbs}})$, dominated by the Cholesky decompositions. With $d=9$ and $N_{\text{gibbs}}=256$, this takes $<1$ms on modern hardware---negligible compared to measurement time.

\paragraph{Leak Probability.}
The leak probability is estimated as the Monte Carlo fraction of posterior samples exceeding the threshold:
\begin{equation}
\widehat{P}(\text{leak} > \theta_{\text{eff}} \mid \Delta) = \frac{1}{N_{\text{keep}}} \sum_{s=1}^{N_{\text{keep}}} \mathbf{1}\!\left[m(\delta^{(s)}) > \theta_{\text{eff}}\right]
\end{equation}

\subsection{Block Bootstrap for Autocorrelation}
\label{app:bootstrap}

\paragraph{Why Standard Bootstrap Fails.}
Timing measurements exhibit autocorrelation: nearby samples are more similar than distant ones due to cache state, CPU frequency scaling, and OS scheduling. Standard bootstrap assumes i.i.d.\ samples, systematically underestimating variance and producing overconfident posteriors.

For example, if consecutive measurements have autocorrelation $\rho = 0.3$, the effective sample size may be 30--50\% lower than the nominal count. Ignoring this inflates the false positive rate---the main failure mode we aim to prevent.

\paragraph{Stream-Based Block Bootstrap.}
We use block bootstrap on the \emph{acquisition stream}---the interleaved sequence of (class, timing) pairs in measurement order---rather than per-class samples. This preserves the true temporal dependence structure.

The algorithm resamples contiguous blocks of length $\hat{b}$ with replacement, reconstructs a bootstrap stream of the same length, splits by class, and computes quantile differences. Over $B = 2000$ iterations, we estimate the covariance matrix $\Sigma_{\text{cal}}$ of $\Delta$.

\paragraph{Block Length Selection.}
We use the Politis-White automatic block length selection algorithm with a class-conditional modification. Standard Politis-White estimates autocorrelation from raw lag-$k$ pairs, but with interleaved sampling, a lag-$k$ pair may cross class boundaries. Instead, we compute autocorrelation using only same-class pairs at each acquisition-stream lag, then take the conservative maximum across classes.

The optimal block length scales as $\hat{b} \propto T^{1/3}$ where $T$ is the stream length. We clamp $\hat{b} \in [10, \min(3\sqrt{T}, T/3)]$ to prevent degenerate blocks. In high-autocorrelation regimes (detected when $\rho_{\text{max}}(10) > 0.3$), we inflate the block length by 50\% for additional conservatism.

\paragraph{Effective Sample Size via IACT.}
We compute the integrated autocorrelation time (IACT) $\hat{\tau}$ using the Geyer initial monotone sequence (IMS) algorithm. This estimates how many nominal samples equal one independent sample:
\begin{equation}
n_{\text{eff}} = \frac{n}{\hat{\tau}}
\end{equation}

For timing analysis, IACT is computed on indicator series $\mathbf{1}[y \leq q_p]$ for each quantile $p$, taking the maximum across quantiles and classes. This captures dependence in the quantile estimators, not just raw timings.

\paragraph{Covariance Scaling.}
The key insight enabling adaptive sampling is that quantile estimator covariance scales as $1/n$. We estimate a ``covariance rate'' during calibration:
\begin{equation}
\Sigma_{\text{rate}} = \hat{\Sigma}_{\text{cal}} \cdot n_{\text{cal}}
\end{equation}
During the adaptive loop with $n$ samples and current IACT $\hat{\tau}$:
\begin{equation}
\Sigma_n = \Sigma_{\text{rate}} \cdot \frac{\hat{\tau}}{\hat{\tau}_{\text{cal}}} \cdot \frac{1}{n}
\end{equation}
This allows cheap posterior updates without re-bootstrapping at each batch.

\subsection{Quality Gates}
\label{app:gates}

Quality gates detect when data quality is insufficient for a confident verdict, forcing an \textsc{Inconclusive} outcome rather than an unreliable Pass or Fail. This is the key mechanism ensuring calibrated verdicts.

\paragraph{Gate 1: Insufficient Information Gain.}
We compute the KL divergence between Gaussian approximations of the prior and posterior:
\begin{multline}
\text{KL}(p_{\text{post}} \| p_{\text{prior}}) = \frac{1}{2}\Big(
\text{tr}(\Lambda_0^{-1}\Lambda_{\text{post}}) \\
+ \mu_{\text{post}}^\top \Lambda_0^{-1} \mu_{\text{post}}
- d + \ln\frac{|\Lambda_0|}{|\Lambda_{\text{post}}|}
\Big)
\end{multline}
where $\Lambda_0 = 2\sigma^2 R$ is the marginal prior covariance for $\nu=4$, and $(\mu_{\text{post}}, \Lambda_{\text{post}})$ are estimated from Gibbs samples.

\textbf{DataTooNoisy}: If KL $< 0.7$ nats after calibration, the posterior has barely moved from the prior---the data is too noisy to distinguish signal from noise.

\textbf{NotLearning}: If the sum of inter-batch KL divergences falls below 0.001 for 5+ consecutive batches, the posterior has stopped updating. Additional samples will not help.

\paragraph{Gate 2: Would Take Too Long.}
We extrapolate time to decision based on current convergence rate. If projected time exceeds 10$\times$ the configured budget, we stop early with \textsc{Inconclusive} rather than waste resources on a test that cannot conclude.

\paragraph{Gate 3: Resource Budget Exceeded.}
If elapsed time exceeds the configured time budget, or total samples exceed the configured maximum, we stop with \textsc{Inconclusive} and report the current posterior probability.

\paragraph{Gate 4: Condition Drift Detected.}
The covariance rate $\Sigma_{\text{rate}}$ is estimated during calibration. If measurement conditions change during the adaptive loop, this estimate becomes invalid. We detect drift by comparing:
\begin{itemize}
\item Variance ratio: $\sigma_{\text{post}}^2 / \sigma_{\text{cal}}^2$ outside $[0.5, 2.0]$
\item Autocorrelation change: $|\rho_{\text{post}}(1) - \rho_{\text{cal}}(1)| > 0.3$
\item Mean drift: $|\mu_{\text{post}} - \mu_{\text{cal}}| / \sigma_{\text{cal}} > 3.0$
\end{itemize}
Any trigger yields \textsc{Inconclusive} with reason \textsc{ConditionsChanged}.

\paragraph{Likelihood Inflation Warning.}
When the posterior mean of $\kappa$ falls below 0.3, the likelihood covariance was effectively scaled up by $>3\times$ for robustness. This indicates potential covariance misestimation and is reported in diagnostics, though it does not block verdicts.

\subsection{Measurement Floor Derivation}
\label{app:floor}

A critical design element is distinguishing what the user \emph{wants} to detect ($\theta_{\text{user}}$) from what the measurement \emph{can} detect ($\theta_{\text{floor}}$).

\paragraph{Floor-Rate Constant.}
During calibration, we compute a floor-rate constant $c_{\text{floor}}$ representing the 95th percentile of the null distribution of $m(\delta)$:
\begin{equation}
c_{\text{floor}} = q_{0.95}\!\left(\max_k |Z_{0,k}|\right), \quad Z_0 \sim \mathcal{N}(0, \Sigma_{\text{rate}})
\end{equation}
estimated via Monte Carlo (50,000 samples). This characterizes the ``noise floor'' of the measurement system.

\paragraph{Dynamic Measurement Floor.}
During the adaptive loop, the statistical floor shrinks with effective sample size:
\begin{equation}
\theta_{\text{floor,stat}}(n) = \frac{c_{\text{floor}}}{\sqrt{n_{\text{eff}}}}
\end{equation}

The timer also imposes a floor based on quantization:
\begin{equation}
\theta_{\text{tick}} = \frac{\text{1 tick (ns)}}{K}
\end{equation}
where $K$ is the batch size (if batching is used). The combined floor is:
\begin{equation}
\theta_{\text{floor}}(n) = \max\!\left(\theta_{\text{floor,stat}}(n), \theta_{\text{tick}}\right)
\end{equation}

\paragraph{Effective Threshold.}
The effective threshold used for inference is:
\begin{equation}
\theta_{\text{eff}} = \max(\theta_{\text{user}}, \theta_{\text{floor}})
\end{equation}

When $\theta_{\text{eff}} > \theta_{\text{user}}$:
\begin{itemize}
\item \textbf{Fail propagates}: Detecting $m(\delta) > \theta_{\text{eff}}$ implies $m(\delta) > \theta_{\text{user}}$, since $\theta_{\text{eff}} \geq \theta_{\text{user}}$.
\item \textbf{Pass does not propagate}: ``No effect above $\theta_{\text{eff}}$'' is compatible with effects in $(\theta_{\text{user}}, \theta_{\text{eff}}]$. We return \textsc{Inconclusive} with reason \textsc{ThresholdElevated}.
\end{itemize}

As sampling continues, $\theta_{\text{floor}}(n)$ decreases. If it drops to $\theta_{\text{user}}$ or below, Pass becomes possible (subject to the posterior). This is why adaptive sampling can eventually reach Pass even when early calibration shows an elevated floor.

\vspace{1em}
\noindent\textit{For complete algorithmic details including pseudocode, numerical stability procedures, and platform-specific considerations, see the specification document in the artifact repository.}


%===============================================================================
% REFERENCES
%===============================================================================
\cleardoublepage
\bibliographystyle{plain}
\bibliography{paper}

\end{document}
